\question[5]
Completa la tabella di confronto dei tre esempi di laboratorio (16, 18, 19): bersaglio, primitiva, difesa chiave e varianti d'architettura supportate sulla vulnbox.
\begin{solution}
\begin{center}\small
\begin{tabular}{@{}lllll@{}}
\toprule
Esempio & Bersaglio & Primitiva & Difesa chiave & Varianti\\
\midrule
16 heap-zero  & puntatore a funzione & overwrite diretto      & NX, bound check   & i386 / x64 / arm64\\
18 heap-two   & puntatore stantio    & use-after-free         & free-and-null     & i386 / x64 / arm64\\
19 heap-three & metadati allocatore  & \lstinline|unlink| / write-what-where & safe unlinking & solo i386\\
\bottomrule
\end{tabular}
\end{center}
Dal dato, al puntatore, ai metadati dell'allocatore: ogni livello colpisce una struttura più profonda ed è una scala di complessità crescente. \lstinline|heap-three| è l'attacco che ha costretto gli allocatori a introdurre i controlli di integrità (safe unlinking): per questo gira solo sulla vecchia dlmalloc 2.7.2 e solo su i386.
\end{solution}
